\documentclass[a4paper,12pt]{jsarticle}
\usepackage[margin=25mm]{geometry}
\usepackage{amsmath,amssymb}
\usepackage{enumitem}

\begin{document}

\begin{center}
{\LARGE 数I・A 共通テスト本番レベル（図形のみ）}
\end{center}

\vspace{1em}

\section*{問題}

\section*{第1問（三角比・図形と計量）}
三角形 \(ABC\) において、
\[
AB=7,\quad AC=9,\quad \angle A=120^\circ
\]
とする。辺 \(BC\) の中点を \(M\) とする。
\begin{enumerate}[label=\arabic*.]
  \item 辺 \(BC\) の長さを求めよ。
  \item 三角形 \(ABC\) の面積を求めよ。
  \item 中線 \(AM\) の長さを求めよ。
  \item 三角形 \(ABC\) の外接円の半径 \(R\) を求めよ。
\end{enumerate}

\section*{第2問（円）}
半径 \(8\) の円 \(O\) において、弦 \(AB\) が
\[
AB=8\sqrt2
\]
を満たしている。
\begin{enumerate}[label=\arabic*.]
  \item 中心角 \(\angle AOB\)（小さい方）を求めよ。
  \item 弧 \(AB\)（小さい方）の長さを求めよ。
  \item 扇形 \(AOB\)（小さい方）の面積を求めよ。
  \item 弓形（弦 \(AB\) と弧 \(AB\) で囲まれる部分）の面積を求めよ。
\end{enumerate}

\section*{第3問（平面図形）}
平行四辺形 \(ABCD\) において、
\[
AB=10,\quad AD=6,\quad \angle DAB=60^\circ
\]
とする。対角線の交点を \(O\) とする。
\begin{enumerate}[label=\arabic*.]
  \item 平行四辺形 \(ABCD\) の面積を求めよ。
  \item 対角線 \(BD\) の長さを求めよ。
  \item 三角形 \(AOD\) の面積を求めよ。
  \item 点 \(A\) から直線 \(BD\) への距離を求めよ。
\end{enumerate}

\section*{第4問（空間図形）}
底面が一辺 \(6\) の正方形、高さ \(9\) の直方体 \(ABCD\text{-}EFGH\) を考える。  
（\(ABCD\) が下底面、\(EFGH\) が上底面、\(AE\perp\) 底面）
\begin{enumerate}[label=\arabic*.]
  \item 空間対角線 \(AG\) の長さを求めよ。
  \item 線分 \(AG\) と底面 \(ABCD\) がなす角 \(\theta\) について、\(\tan\theta\) を求めよ。
  \item 三角形 \(AEG\) の面積を求めよ。
  \item 点 \(A\) から平面 \(BDG\) までの距離を求めよ。
\end{enumerate}

\end{document}
